\itemhead{第1章 小结与复习 · 第3题（教材p.28）}{XJ-B1-C01-SU-P028-Y1232}
\label{XJ-B1-C01-SU-P028-Y1232}
{\footnotesize 湘教版教材p.28 / PDF第35页 / 第1章 小结与复习 / 第3题}
\field{原题}{集合是数学的基本概念，数学概念之间的关系，需要用逻辑用语来表达。学习常用逻辑用语可以更深刻地理解数学内容，更严密地进行数学推理以及更准确地表达数学思想。}
\field{参考答案}{这是概念联系的总结项。例如 $A\subseteq B$ 正是“$x\in A$”推出“$x\in B$”；将包含关系与充分条件联系，能检查推理方向，不另设唯一数值答案。}
{\small \field{关联知识}{L02 充分、必要、充要条件（教材帮p.30、31 / PDF39、40）}
\field{方法与易错点}{分别检查两个方向；用解集包含或具体反例判定，不凭句子位置猜方向。 易错：把“必要”箭头反写；只有一个方向证据便断言“不必要”；忽略给定论域。}
\field{答案依据}{答案册未提供本栏目的独立对应答案；依据原题、教材概念补写参考回应。开放性总结不设唯一数值答案。}}
